% IEEE Conference Paper — Combined Attacks on ML
% Expanded version (~12 pages)
\documentclass[conference]{IEEEtran}

% ─── Packages ─────────────────────────────────────────────────────────────────
\usepackage{cite}
\usepackage{amsmath,amssymb,amsfonts}
\usepackage{algorithmic}
\usepackage{graphicx}
\usepackage{textcomp}
\usepackage{xcolor}
\usepackage{booktabs}
\usepackage{multirow}
\usepackage{subcaption}
\usepackage{hyperref}
\usepackage{algorithm}
\usepackage{algpseudocode}
\usepackage{enumitem}

\def\BibTeX{{\rm B\kern-.05em{\sc i\kern-.025em b}\kern-.08em
    T\kern-.1667em\lower.7ex\hbox{E}\kern-.125emX}}

\begin{document}

% ═══════════════════════════════════════════════════════════════════════════════
% TITLE & AUTHORS
% ═══════════════════════════════════════════════════════════════════════════════

\title{Exploiting the Robustness-Privacy Tradeoff: Combined Evasion and Membership Inference Attacks on Adversarially Trained Medical Image Classifiers}

\author{
    \IEEEauthorblockN{Author Name\textsuperscript{1}}
    \IEEEauthorblockA{
        \textit{Department of Computer Science} \\
        \textit{Institute Name} \\
        City, Country \\
        email@institute.edu
    }
    \and
    \IEEEauthorblockN{Co-Author Name\textsuperscript{2}}
    \IEEEauthorblockA{
        \textit{Department of Computer Science} \\
        \textit{Institute Name} \\
        City, Country \\
        email@institute.edu
    }
}

\maketitle

% ═══════════════════════════════════════════════════════════════════════════════
% ABSTRACT
% ═══════════════════════════════════════════════════════════════════════════════

\begin{abstract}
Deep learning models deployed in safety-critical medical imaging applications face two fundamental classes of adversarial threats: \textit{evasion attacks}, which craft adversarial inputs to cause misclassification at inference time, and \textit{membership inference attacks} (MIA), which determine whether a specific data record was used in the model's training set, thereby violating patient privacy. Adversarial training, the predominant defense against evasion attacks, has been widely adopted to improve model robustness. However, we demonstrate that this defense mechanism introduces a critical vulnerability: \textit{robust overfitting}, wherein the model memorizes adversarial patterns specific to its training data. We exploit this phenomenon through a novel \textbf{fluctuation-based membership inference attack} that measures the differential stability of model outputs under adversarial perturbation---without requiring any knowledge of the victim's training procedure. Our experiments on a DenseNet-121 model trained on the NIH Chest X-ray dataset show that while adversarial training ($\varepsilon = 0.02$) successfully reduces evasion attack success, it simultaneously elevates MIA accuracy from $\sim$51\% (random guessing) on a standard baseline model to $\sim$70\% on the adversarially trained model. These findings reveal a fundamental \textit{robustness--privacy tradeoff} in medical AI systems and underscore the need for defenses that jointly address both threat vectors.
\end{abstract}

\begin{IEEEkeywords}
adversarial machine learning, membership inference attack, evasion attack, adversarial training, robust overfitting, medical image classification, privacy
\end{IEEEkeywords}

% ═══════════════════════════════════════════════════════════════════════════════
% I. INTRODUCTION
% ═══════════════════════════════════════════════════════════════════════════════

\section{Introduction}

The adoption of deep learning in medical image analysis has progressed rapidly in recent years, with convolutional neural networks achieving radiologist-level performance on tasks such as pneumonia detection from chest X-rays~\cite{rajpurkar2017chexnet}, diabetic retinopathy screening~\cite{gulshan2016development}, and skin lesion classification~\cite{esteva2017dermatologist}. These models are increasingly being considered for clinical deployment, where they assist or augment human decision-making in high-stakes diagnostic settings.

However, deploying machine learning models in healthcare introduces a dual set of concerns that are often addressed in isolation. On the \textit{security} front, models are vulnerable to \textit{evasion attacks}---carefully crafted perturbations to input images that can cause a trained classifier to produce incorrect predictions. A manipulated X-ray that causes a model to miss a pneumothorax, for instance, could have life-threatening consequences. On the \textit{privacy} front, the very act of training on sensitive patient data creates the risk of \textit{membership inference attacks} (MIA), which can reveal whether a specific patient's medical image was part of the training dataset. Such a disclosure constitutes a violation of patient confidentiality, potentially running afoul of regulations such as HIPAA in the United States and GDPR in the European Union.

These two threats are typically studied independently: the adversarial robustness community develops defenses against evasion attacks, while the privacy community develops techniques to protect against membership inference. In this paper, we demonstrate that the dominant defense against evasion attacks---\textit{adversarial training}---while effective at its primary goal, creates a previously underexplored vulnerability to membership inference attacks. The mechanism is \textit{robust overfitting}~\cite{rice2020overfitting}: an adversarially trained model develops strong defenses against adversarial perturbations on its training data, but this robustness fails to generalize to unseen data. This creates a detectable behavioral asymmetry that an attacker can exploit to infer membership.

% DIAGRAM PLACEHOLDER 1
\begin{figure}[t]
    \centering
    \includegraphics[width=\columnwidth]{images/img1.png}
    % \fbox{\parbox{0.92\columnwidth}{\centering\vspace{3.5cm}
    % \textbf{[DIAGRAM: System Overview]}\\[4pt]
    % A high-level overview showing: (1) A medical institution trains a DenseNet-121 model on patient X-rays and deploys it as a diagnostic API. (2) An evasion attacker crafts adversarial X-rays to cause misdiagnosis. (3) A membership inference attacker queries the API to determine if a specific patient's X-ray was in the training set. (4) The institution applies adversarial training as a defense against (2), but this inadvertently enables (3).\\
    % \textit{Suggested style:} Block diagram with hospital/model/attacker icons and arrows showing the two attack vectors.
    % \vspace{0.5cm}}}
    \caption{Overview of the combined attack scenario. A medical institution deploys a model defended by adversarial training. While this successfully mitigates evasion attacks, it inadvertently introduces a membership inference vulnerability through robust overfitting.}
    \label{fig:system_overview}
\end{figure}

\subsection{Contributions}

In this paper, we make the following contributions:

\begin{enumerate}[leftmargin=*]
    \item We demonstrate the existence of a \textit{robustness--privacy tradeoff} on a real-world medical imaging task (NIH Chest X-ray multi-label classification with DenseNet-121), showing that adversarial training elevates MIA accuracy from $\sim$51\% to $\sim$70\%.
    
    \item We propose a novel \textbf{fluctuation-based MIA} that trains the attack classifier on fluctuation features---$|\mathbf{p}_{\text{clean}} - \mathbf{p}_{\text{adv}}|$---rather than raw confidence scores. This attack operates under a strict black-box threat model and requires \textit{no knowledge} of the victim's training procedure.
    
    \item We introduce a \textbf{variance-enhanced feature} (bootstrapped variance of max-confidence scores) that provides complementary discriminative information. This feature captures model-level behavioral statistics and can serve as a general-purpose enhancement for any MIA variant, including the standard shadow-model approach~\cite{shokri2017membership}, by injecting richer context into the per-sample feature vector.
\end{enumerate}

The remainder of this paper is organized as follows. Section~\ref{sec:related_work} surveys related work on adversarial attacks, membership inference, and the privacy implications of robust training. Section~\ref{sec:threat_model} describes the threat model and experimental setup. Section~\ref{sec:methodology} details the attack methodologies. Section~\ref{sec:results} presents the experimental results, and Section~\ref{sec:conclusion} concludes.


% ═══════════════════════════════════════════════════════════════════════════════
% II. RELATED WORK
% ═══════════════════════════════════════════════════════════════════════════════

\section{Related Work}
\label{sec:related_work}

\subsection{Adversarial Evasion Attacks}

The vulnerability of deep neural networks to adversarial examples was first identified by Szegedy et al.~\cite{szegedy2014intriguing}, who observed that imperceptible perturbations to input images could cause state-of-the-art classifiers to produce incorrect predictions with high confidence. Goodfellow et al.~\cite{goodfellow2015explaining} provided a theoretical explanation rooted in the local linearity of deep networks and proposed the Fast Gradient Sign Method (FGSM), a single-step attack that computes the perturbation as the sign of the input-space loss gradient. Madry et al.~\cite{madry2018towards} strengthened this with Projected Gradient Descent (PGD), an iterative multi-step variant that provides stronger attacks by solving a constrained optimization problem over multiple gradient steps. Carlini and Wagner~\cite{carlini2017towards} proposed optimization-based attacks that can defeat many proposed defenses by directly optimizing for misclassification under minimal perturbation.

In the medical imaging domain, Finlayson et al.~\cite{finlayson2019adversarial} demonstrated that adversarial attacks can manipulate medical AI systems to commit insurance fraud by altering diagnostic outputs, highlighting the real-world stakes of evasion attacks in healthcare. Paschali et al.~\cite{paschali2018generalizability} studied the generalizability of adversarial attacks across medical imaging modalities, showing that attacks transfer across architectures and imaging tasks.

\subsection{Defenses Against Evasion Attacks}

Adversarial training, first proposed by Goodfellow et al.~\cite{goodfellow2015explaining} and formalized by Madry et al.~\cite{madry2018towards}, remains the most widely adopted and empirically effective defense. The key idea is to augment the training data with adversarial examples generated on the fly, so the model learns to produce correct predictions under perturbation. The training objective becomes a min-max optimization: the inner maximization finds the worst-case perturbation, while the outer minimization trains the model to withstand it.

However, Rice et al.~\cite{rice2020overfitting} identified a fundamental failure mode of adversarial training called \textit{robust overfitting}: as training progresses, the model's adversarial robustness on the training set continues to improve, but its robustness on the test set stagnates or even degrades. This occurs because the model memorizes specific adversarial patterns for its training examples rather than learning generalizable robustness. Several mitigation strategies have been proposed, including early stopping on adversarial validation loss~\cite{rice2020overfitting}, adversarial weight perturbation~\cite{wu2020adversarial}, and stochastic weight averaging~\cite{chen2021robust}. Despite these efforts, robust overfitting remains a pervasive challenge.

Other defense strategies include input transformation~\cite{guo2018countering}, defensive distillation~\cite{papernot2016distillation}, and certified defenses based on randomized smoothing~\cite{cohen2019certified}. However, adversarial training remains the benchmark approach and is the focus of this work.

\subsection{Membership Inference Attacks}

Membership Inference Attacks (MIA) were introduced by Shokri et al.~\cite{shokri2017membership}, who proposed a shadow-model-based approach: the attacker trains multiple surrogate models to mimic the victim's behavior, then uses the behavioral differences between members and non-members of the shadow models' training sets to train a binary attack classifier. When applied to the victim's outputs, this classifier can predict whether a given data point was used for training.

Salem et al.~\cite{salem2019ml} showed that the attack can be significantly simplified---fewer shadow models, relaxed architectural assumptions, and even threshold-based attacks can be effective when the target model exhibits sufficient overfitting. Yeom et al.~\cite{yeom2018privacy} established theoretical connections between MIA success and the model's generalization gap, showing that the probability of a successful attack is bounded by the difference between training and test accuracy. Song and Mittal~\cite{song2021systematic} provided a systematic evaluation of MIA across different model architectures and training procedures, demonstrating that factors such as model complexity, training duration, and regularization all influence MIA vulnerability.

Label-only attacks~\cite{choquette2021label} and metric-based attacks~\cite{carlini2022membership} have further expanded the MIA toolkit, showing that even limited model access can leak membership information.

\subsection{Privacy Risks of Adversarial Training}

The intersection of adversarial robustness and privacy has received growing attention. Song et al.~\cite{song2019privacy} were the first to identify that adversarial training \textit{increases} a model's vulnerability to membership inference, demonstrating this on CIFAR-10 with PGD-adversarially trained models. Their key observation was that robust models exhibit larger generalization gaps in adversarial accuracy (i.e., robust overfitting), which amplifies the membership signal.

We build on this line of work by: (i)~extending the study to a real-world multi-label medical imaging task (NIH Chest X-rays) with direct clinical implications, and (ii)~proposing a novel fluctuation-based MIA that directly uses adversarial fluctuation (the change in output under perturbation) as the membership feature, which we show is the key to unlocking the robust overfitting signal.


% ═══════════════════════════════════════════════════════════════════════════════
% III. THREAT MODEL
% ═══════════════════════════════════════════════════════════════════════════════

\section{Threat Model}
\label{sec:threat_model}

We consider a realistic scenario in which a medical institution trains and deploys a deep learning model for automated chest X-ray diagnosis. An adversary seeks to exploit this deployed model through two distinct but related attack vectors.

\subsection{The Victim Model}

\subsubsection{Architecture}
The victim model is a DenseNet-121~\cite{huang2017densely} pre-trained on ImageNet~\cite{deng2009imagenet} and fine-tuned on the NIH Chest X-ray dataset~\cite{wang2017chestxray}. DenseNet-121 was chosen because it is the standard architecture for chest X-ray classification, following CheXNet~\cite{rajpurkar2017chexnet}. The final classification head replaces the default 1000-class ImageNet classifier with a custom multi-label head:
\begin{equation}
    \text{Head}: \text{Linear}(1024, 512) \rightarrow \text{ReLU} \rightarrow \text{Dropout}(0.3) \rightarrow \text{Linear}(512, 15)
\end{equation}
This head produces 15 logits corresponding to 14 pathology classes (Atelectasis, Consolidation, Infiltration, Pneumothorax, Edema, Emphysema, Fibrosis, Effusion, Pneumonia, Pleural Thickening, Cardiomegaly, Nodule, Mass, Hernia) and one ``No Finding'' class. A sigmoid activation is applied to each logit independently, producing multi-label confidence scores in $[0, 1]^{15}$.

\subsubsection{Training Loss}
The model is trained with Binary Cross-Entropy with Logits (BCEWithLogitsLoss) with class-balanced positive weights:
\begin{equation}
    \mathcal{L} = -\frac{1}{C} \sum_{c=1}^{C} \Big[ w_c^+ \cdot y_c \log(\sigma(z_c)) + (1 - y_c) \log(1 - \sigma(z_c)) \Big]
\end{equation}
where $C = 15$ is the number of classes, $z_c$ is the logit for class $c$, $\sigma$ is the sigmoid function, and $w_c^+ = \min\left(\frac{N_c^-}{N_c^+}, 10\right)$ is the positive weight for class $c$, clipped to prevent extreme imbalance (necessary because rare conditions like Hernia have very few positive examples).

\subsubsection{Two Victim Variants}
To study the robustness--privacy tradeoff, we train two variants of the victim model:

\begin{enumerate}[leftmargin=*]
    \item \textbf{Baseline Model.} Trained using standard best practices: AdamW optimizer ($lr = 10^{-4}$, weight decay $= 10^{-4}$), linear warmup for 3 epochs followed by cosine annealing, data augmentation (random resized crop, horizontal flip, rotation, color jitter), Dropout(0.3), and early stopping on validation AUC with patience 7. This represents a well-regularized, production-quality model.
    
    \item \textbf{Adversarial Model.} Trained with FGSM adversarial training (detailed in Section~\ref{sec:adv_training_detail}) using $\varepsilon = 0.02$, $\alpha = 0.5$, for 35 epochs \textit{without} data augmentation or early stopping. Augmentation is disabled to isolate the effect of adversarial training; early stopping is disabled to allow robust overfitting to fully manifest.
\end{enumerate}

\subsection{Dataset and Partition}

The NIH Chest X-ray dataset contains 112,120 frontal-view chest X-ray images from 30,805 unique patients, with multi-label annotations for 15 classes. We partition this dataset as follows:

\begin{itemize}[leftmargin=*]
    \item \textbf{Members} ($\sim$78,708 images): Used to train the victim model. A 15\% internal validation split (11,806 images) is reserved for hyperparameter tuning, with the remaining 66,902 images used for training.
    \item \textbf{Non-members} ($\sim$33,412 images): Completely withheld from the victim's training process. Used by the attacker for evaluation and as part of the unlabelled pool.
\end{itemize}

All images are resized to $224 \times 224$ pixels and normalized using ImageNet statistics ($\mu = [0.485, 0.456, 0.406]$, $\sigma = [0.229, 0.224, 0.225]$).

\subsection{Attacker's Knowledge and Capabilities}

\subsubsection{Evasion Attacker (White-Box)}
The evasion attacker has \textit{full white-box access} to the victim model, including its architecture, trained weights, and gradients. This is the strongest possible attacker model and serves as an upper bound on evasion attack effectiveness. The attacker applies FGSM (Eq.~\ref{eq:fgsm}) to craft adversarial X-ray images under an $\ell_\infty$ perturbation budget $\varepsilon$.

\subsubsection{Membership Inference Attacker (Black-Box)}
The MIA attacker operates under a strict \textit{black-box} threat model with minimal assumptions:
\begin{itemize}[leftmargin=*]
    \item \textbf{Query access only.} The attacker can submit image file paths to the victim's prediction API and receive back a 15-dimensional sigmoid confidence vector. The attacker cannot inspect the model's architecture, weights, gradients, or internal activations.
    
    \item \textbf{Unlabelled data pool.} The attacker possesses a large pool of images drawn from the same distribution as the victim's training data, but does \textit{not} know which specific images are members or non-members. We simulate this by shuffling 20,000 member images and 20,000 non-member images into a single unlabelled pool of 40,000 images.
    
    \item \textbf{No knowledge of the victim's training procedure.} Critically, the attacker does \textit{not} know whether the victim was trained using standard training, adversarial training, or any other specific method. The attacker also does not know the victim's architecture, hyperparameters, or regularization strategy. The proposed fluctuation-based attack is designed as a general-purpose probing technique that does not require such knowledge.
    
    \item \textbf{Unlimited queries.} The attacker can query the victim API as many times as needed. This is realistic for a deployed medical AI API without rate limiting.
\end{itemize}

The attacker's goal is to determine, for each image in a held-out evaluation set of 5,000 members and 5,000 non-members, whether the image was part of the victim's training set.


% ═══════════════════════════════════════════════════════════════════════════════
% IV. METHODOLOGY
% ═══════════════════════════════════════════════════════════════════════════════

\section{Methodology}
\label{sec:methodology}

This section describes the four core components of our work: (1)~the FGSM evasion attack and how adversarial training defends against it, (2)~the standard shadow-model-based MIA framework that our approach builds upon, (3)~the variance-enhanced feature as a general-purpose improvement, and (4)~our proposed fluctuation-based MIA that exploits robust overfitting.

% ─── IV-A. EVASION ATTACK ─────────────────────────────────────────────────────

\subsection{Evasion Attack: Fast Gradient Sign Method (FGSM)}
\label{sec:evasion_attack}

The Fast Gradient Sign Method~\cite{goodfellow2015explaining} exploits the local linearity of deep neural networks to craft adversarial examples in a single forward-backward pass. The key insight is that the loss function $\mathcal{L}$ of a neural network is approximately linear in the input space for small perturbations. Under this linear approximation, the perturbation that maximally increases the loss under an $\ell_\infty$ constraint is the sign of the gradient.

% DIAGRAM PLACEHOLDER 2
\begin{figure}[t]
    \centering
    % \includegraphics[width=\columnwidth]{figures/fgsm_attack.png}
    \fbox{\parbox{0.92\columnwidth}{\centering\vspace{3cm}
    \textbf{[DIAGRAM: FGSM Attack Visualization]}\\[4pt]
    Show 3 images side by side: (a) Original clean chest X-ray with correct prediction ``Pneumonia: 0.87'', (b) The gradient-sign perturbation pattern $\varepsilon \cdot \text{sign}(\nabla_x \mathcal{L})$ (visualized as a noise pattern), (c) The resulting adversarial image (visually identical to original) with incorrect prediction ``Pneumonia: 0.12''.\\[4pt]
    Below the images, show the equation: $x_{adv} = x + \varepsilon \cdot \text{sign}(\nabla_x \mathcal{L})$.\\
    \textit{Suggested style:} Three chest X-ray images with arrows and ``+'' sign between them. Annotate with confidence scores.
    \vspace{0.5cm}}}
    \caption{FGSM evasion attack on a chest X-ray. A visually imperceptible perturbation (magnified in center) is added to the clean image, causing the model to dramatically reduce its confidence in the correct diagnosis.}
    \label{fig:fgsm_attack}
\end{figure}

Concretely, given a trained model $f_\theta$, a clean input image $x$ (a chest X-ray), and its corresponding multi-hot label vector $y \in \{0, 1\}^{15}$, the FGSM attack proceeds as follows:

\begin{enumerate}[leftmargin=*]
    \item \textbf{Forward pass.} Compute the model's output logits $z = f_\theta(x)$ and the loss $\mathcal{L}(z, y)$ using Binary Cross-Entropy with Logits.
    
    \item \textbf{Backward pass.} Compute the gradient of the loss with respect to the \textit{input image}: $g = \nabla_x \mathcal{L}(f_\theta(x), y)$. This gradient tells us, for each pixel, which direction of change would most increase the loss (i.e., most confuse the model).
    
    \item \textbf{Perturbation.} Construct the adversarial image by stepping in the direction of the gradient sign:
    \begin{equation}
        x_{adv} = \text{clamp}\Big(x + \varepsilon \cdot \text{sign}(g), \; x_{\min}, \; x_{\max}\Big)
        \label{eq:fgsm}
    \end{equation}
    where $\varepsilon$ controls the perturbation magnitude and the clamping ensures the result stays within valid pixel ranges. In our implementation, images are normalized to ImageNet statistics, so we clamp to $[-3.0, 3.0]$ in the normalized space.
\end{enumerate}

The perturbation magnitude $\varepsilon$ controls the tradeoff between attack strength and imperceptibility. We evaluate across $\varepsilon \in \{0.001, 0.005, 0.01, 0.02, 0.05, 0.1\}$ to characterize the full attack curve. At $\varepsilon = 0.01$ in normalized pixel space, the perturbation is imperceptible to the human eye but can significantly degrade model performance.

% ─── IV-B. ADVERSARIAL TRAINING ───────────────────────────────────────────────

\subsection{Adversarial Training: The Evasion Defense}
\label{sec:adv_training_detail}

Adversarial training~\cite{goodfellow2015explaining, madry2018towards} is the process of augmenting the training procedure with adversarial examples so the model learns to produce correct predictions under perturbation. It is widely considered the most effective defense against evasion attacks. In this section, we describe our implementation in detail, as the specific training procedure is critical to understanding why it creates a membership inference vulnerability.

% DIAGRAM PLACEHOLDER 3
\begin{figure}[t]
    \centering
    % \includegraphics[width=\columnwidth]{figures/adv_training_loop.png}
    \fbox{\parbox{0.92\columnwidth}{\centering\vspace{4cm}
    \textbf{[DIAGRAM: Adversarial Training Loop]}\\[4pt]
    A flowchart showing one training iteration:\\
    (1) Input batch of clean images $\{x_i\}$.\\
    (2) Forward pass through model $\rightarrow$ clean logits.\\
    (3) Compute $\mathcal{L}_{\text{clean}}$ and backprop to get $\nabla_x \mathcal{L}$.\\
    (4) Generate adversarial images: $x_{adv} = x + \varepsilon \cdot \text{sign}(\nabla_x \mathcal{L})$.\\
    (5) Forward pass adversarial images $\rightarrow$ adversarial logits.\\
    (6) Compute combined loss: $\alpha \mathcal{L}_{\text{clean}} + (1-\alpha) \mathcal{L}_{\text{adv}}$.\\
    (7) Backprop combined loss and update model weights $\theta$.\\[4pt]
    \textit{Suggested style:} Vertical flowchart with model block in center, two forward-pass branches (clean and adversarial), merging at the combined loss.
    \vspace{0.5cm}}}
    \caption{One iteration of the FGSM adversarial training loop. The model first generates adversarial examples using its own gradients, then updates its weights on a combined clean + adversarial loss.}
    \label{fig:adv_training_loop}
\end{figure}

\subsubsection{Training Loop}

For each mini-batch of training images $\{(x_i, y_i)\}_{i=1}^{B}$ during adversarial training, the following steps are performed:

\textbf{Step 1: Generate adversarial examples.} The current model is used to compute the input-space gradient of the loss:
\begin{equation}
    g_i = \nabla_{x_i} \mathcal{L}\big(f_\theta(x_i),\; y_i\big)
\end{equation}
The adversarial image is then constructed via FGSM:
\begin{equation}
    x_{i,adv} = \text{clamp}\big(x_i + \varepsilon \cdot \text{sign}(g_i),\; -3.0,\; 3.0\big)
\end{equation}
Note that this step uses the model's \textit{current} weights $\theta$ at each iteration, meaning the adversarial examples evolve as the model trains. Early in training, the perturbations are less targeted; later, they become increasingly specific to the model's learned decision boundary.

\textbf{Step 2: Compute combined loss.} The model is evaluated on \textit{both} the clean and adversarial images, and their losses are combined:
\begin{equation}
    \mathcal{L}_{\text{total}} = \alpha \cdot \mathcal{L}_{\text{clean}} + (1 - \alpha) \cdot \mathcal{L}_{\text{adv}}
    \label{eq:adv_training}
\end{equation}
where $\mathcal{L}_{\text{clean}} = \frac{1}{B}\sum_i \mathcal{L}(f_\theta(x_i), y_i)$ and $\mathcal{L}_{\text{adv}} = \frac{1}{B}\sum_i \mathcal{L}(f_\theta(x_{i,adv}), y_i)$. We use $\alpha = 0.5$, giving equal weight to clean accuracy and adversarial robustness.

\textbf{Step 3: Update weights.} The gradients of $\mathcal{L}_{\text{total}}$ with respect to the model parameters $\theta$ are computed and used to update the weights via AdamW with gradient clipping (max norm = 5.0).

\subsubsection{Hyperparameters}
The adversarial model is trained with the following configuration:
\begin{itemize}[leftmargin=*]
    \item \textbf{Perturbation magnitude:} $\varepsilon = 0.02$ (in normalized pixel space).
    \item \textbf{Loss balance:} $\alpha = 0.5$ (equal clean and adversarial weight).
    \item \textbf{Optimizer:} AdamW ($lr = 10^{-4}$, weight decay $= 10^{-4}$).
    \item \textbf{LR schedule:} Linear warmup for 3 epochs (start factor 0.1) followed by cosine annealing.
    \item \textbf{Epochs:} 35 (no early stopping).
    \item \textbf{Data augmentation:} Disabled (to isolate adversarial training effects).
    \item \textbf{Mixed precision:} Enabled (FP16 via torch.amp on CUDA).
    \item \textbf{Regularization:} Dropout(0.3) in the classification head, weight decay $10^{-4}$.
\end{itemize}

\subsubsection{Robust Overfitting: The Side Effect}
\label{sec:robust_overfitting}

A critical phenomenon observed during adversarial training is \textit{robust overfitting}~\cite{rice2020overfitting}. As training progresses, the model's \textit{clean accuracy} on both training and validation data converges (training: 99.99\%, validation: 92.32\%, a modest 7.67\% gap). However, the model's \textit{adversarial robustness} diverges sharply:

\begin{itemize}[leftmargin=*]
    \item On \textbf{members} (training data): The model becomes highly robust to adversarial perturbations, maintaining high accuracy under FGSM attack. Its confidence scores barely fluctuate when noise is added.
    
    \item On \textbf{non-members} (unseen data): The model's robustness is significantly weaker. Under the same FGSM perturbation, its confidence scores fluctuate dramatically.
\end{itemize}

We quantify this asymmetry by computing the \textit{average fluctuation} per class:
\begin{equation}
    \text{Fluct}_{c}^{(\text{set})} = \frac{1}{|\mathcal{S}|} \sum_{x \in \mathcal{S}} \left| \sigma(f_\theta(x))_c - \sigma(f_\theta(x_{adv}))_c \right|
    \label{eq:fluctuation_def}
\end{equation}
where $\mathcal{S}$ is either the member or non-member set, and $\sigma$ is the sigmoid function. Our experiments reveal that the average fluctuation on non-members (7.41\%) is nearly \textit{three times} that on members (2.67\%), creating a strong signal for membership inference. This insight forms the foundation of our proposed attack.

Importantly, the clean confidence scores do not exhibit this gap. Both members and non-members produce similar average max-confidence scores ($\sim$0.89) under the adversarially trained model. A conventional MIA approach based solely on clean confidence patterns would therefore fail to detect any membership signal, while our fluctuation-based approach---which actively probes the model's adversarial response---succeeds at $\sim$70\%.

% ─── IV-C. STANDARD MIA ──────────────────────────────────────────────────────

\subsection{Standard Shadow-Model MIA Framework}
\label{sec:standard_mia}

Our fluctuation-based attack builds upon the shadow-model MIA framework introduced by Shokri et al.~\cite{shokri2017membership}. This four-stage pipeline uses \textit{shadow models} to learn the behavioral differences between a model's treatment of members versus non-members. We describe each stage in detail, as it forms the foundation that our approach (Section~\ref{sec:fluctuation_mia}) extends and modifies.

% DIAGRAM PLACEHOLDER 4
\begin{figure*}[t]
    \centering
    \includegraphics[width=\textwidth]{images/img4.png}
    % \fbox{\parbox{0.95\textwidth}{\centering\vspace{4cm}
    % \textbf{[DIAGRAM: Standard MIA Pipeline (Full-Width)]}\\[4pt]
    % A left-to-right pipeline diagram with 4 stages:\\
    % \textbf{Stage 1 (Shadow Training):} Attacker's pool of 40K unlabelled images $\rightarrow$ Sample 10K subset $\rightarrow$ Query Victim API to get pseudo-labels $\rightarrow$ Train shadow model on (images, pseudo-labels). Show 8 shadow models being trained in parallel with different architectures (ResNet-18, MobileNet, EfficientNet, DenseNet-121, ShuffleNet, cycling).\\
    % \textbf{Stage 2 (Attack Dataset):} For each shadow model: query it on its own training data (``members'') and on data it didn't train on (``non-members'') $\rightarrow$ collect 15-dim confidence vectors $\rightarrow$ label as member=1 / non-member=0.\\
    % \textbf{Stage 3 (Attack Model):} Train a Gradient Boosting classifier on the labeled confidence vectors.\\
    % \textbf{Stage 4 (Inference):} New image $\rightarrow$ query victim API $\rightarrow$ get 15-dim confidence $\rightarrow$ feed to attack model $\rightarrow$ predict ``member'' or ``non-member''.\\[4pt]
    % \textit{Suggested style:} Horizontal flow with distinct colored boxes for each stage and small model/data icons.
    % \vspace{0.5cm}}}
    \caption{The standard shadow-model MIA pipeline. The attacker trains shadow models to mimic the victim, builds an attack dataset from the shadow models' differential behavior on members vs.\ non-members, and trains an attack classifier to predict membership from the victim's confidence output.}
    \label{fig:standard_mia}
\end{figure*}

\subsubsection{Stage 1: Shadow Model Training}

The attacker does not have direct access to the victim's training data or its membership labels. Instead, the attacker possesses a large \textit{unlabelled pool} of images from the same distribution. This pool is a shuffled mixture of 20,000 member images and 20,000 non-member images---the attacker cannot distinguish between them.

For each of the $K = 8$ shadow models, the attacker:

\begin{enumerate}[leftmargin=*]
    \item \textbf{Samples a training subset.} Randomly selects $m = 10{,}000$ images from the 40,000-image pool (without replacement within a single shadow model, but with overlap across shadow models).
    
    \item \textbf{Obtains pseudo-labels from the victim.} Queries the victim API on the selected images to obtain the victim's 15-dimensional sigmoid confidence vectors. These are thresholded at 0.5 to produce multi-hot pseudo-labels: $\hat{y}_i = \mathbb{1}[f_V(x_i) > 0.5]$. The use of pseudo-labels is a deliberate design choice that serves a dual purpose: (a)~it enables training without access to ground-truth annotations, and more importantly, (b)~it forces each shadow model to learn the \textit{victim's specific decision boundary} rather than the true data distribution. By training on the victim's own predictions, the shadow models internalize the victim's biases, error patterns, and confidence calibration---faithfully mimicking how the victim model treats its own members versus unseen data. This mimicry is essential for the attack, because the membership signal lies in the \textit{model-specific} behavioral gap, not in the data labels themselves.
    
    \item \textbf{Trains the shadow model.} A PyTorch CNN (pretrained on ImageNet) is fine-tuned on the $(x_i, \hat{y}_i)$ pairs using BCEWithLogitsLoss with the same class-balanced positive weights as the victim. To ensure diversity, the 8 shadow models cycle through 5 different architectures: ResNet-18, MobileNetV3-Small, EfficientNet-B0, DenseNet-121, and ShuffleNetV2-x1.0.
\end{enumerate}

The resulting shadow models approximately mimic the victim's behavior: they produce similar confidence distributions on similar inputs, without requiring access to the victim's architecture or weights.

\subsubsection{Stage 2: Attack Dataset Construction}

The attack dataset captures the behavioral difference between how a model treats its training data (members) versus unseen data (non-members). For each shadow model $s_j$ ($j = 1, \ldots, K$):

\begin{enumerate}[leftmargin=*]
    \item \textbf{Identify members and non-members.} The shadow model's training subset is known to the attacker (since they constructed it). Images in this subset are labeled as members (1); images from the pool \textit{not} in this subset are labeled as non-members (0).
    
    \item \textbf{Query the shadow model.} The shadow model is queried on a balanced sample of members and non-members, producing 15-dimensional confidence vectors for each image.
    
    \item \textbf{Construct feature vectors.} Each sample in the attack dataset consists of:
    \begin{itemize}
        \item The 15 class-wise confidence scores: $\big(\sigma(z_1), \ldots, \sigma(z_{15})\big)$.
        \item The maximum confidence score: $\max_c \sigma(z_c)$.
        \item The class index of the maximum confidence: $\arg\max_c \sigma(z_c)$.
        \item The membership label: 1 (member) or 0 (non-member).
    \end{itemize}
\end{enumerate}

The combined dataset across all 8 shadow models provides a large, diverse training set for the attack classifier.

\subsubsection{Stage 3: Attack Model Training}

A supervised binary classifier---the \textit{attack model}---is trained on the attack dataset to predict membership from the 17-dimensional feature vector (15 confidence scores + max confidence + max class). Common choices for the attack model include Gradient Boosting, Random Forest, MLP, and Logistic Regression.

\subsubsection{Stage 4: Inference}

To attack a new image $x$:
\begin{enumerate}[leftmargin=*]
    \item Query the victim API: $\mathbf{p} = f_V(x) \in [0, 1]^{15}$.
    \item Construct the 17-dim feature vector: $(\mathbf{p}, \max(\mathbf{p}), \arg\max(\mathbf{p}))$.
    \item Feed to the trained attack model: $\hat{m} = \text{AttackModel}(\text{features})$.
    \item Output: $\hat{m} = 1$ (predicted member) or $\hat{m} = 0$ (predicted non-member).
\end{enumerate}

\subsubsection{Motivation for the Fluctuation-Based Approach}

Our approach is motivated by the observation that adversarially trained models may produce nearly identical clean confidence distributions for members and non-members, yet their responses to adversarial perturbation differ significantly. Instead of examining clean outputs, we \textit{actively probe} the model with adversarial perturbations and measure how much its outputs change, revealing the robust overfitting signal.

% ─── IV-D. VARIANCE-ENHANCED MIA ──────────────────────────────────────────────

\subsection{Variance-Enhanced MIA: A General-Purpose Improvement}
\label{sec:variance_mia}

Before presenting our fluctuation-based approach, we describe a general-purpose feature enhancement that can be applied to \textit{any} MIA variant: the \textbf{variance-of-max} feature. This technique, inspired by the data-augmentation-based defenses studied in~\cite{salem2019ml}, injects model-level behavioral statistics into the per-sample feature vector, providing complementary information that improves attack accuracy.

\subsubsection{Motivation}

The standard MIA feature vector (Section~\ref{sec:standard_mia}) consists entirely of \textit{per-sample} features: the model's confidence scores for a single query image. However, membership inference is fundamentally about the relationship between a \textit{sample} and a \textit{model}. By augmenting the per-sample features with a statistic that characterizes the model's \textit{overall} behavior, we provide the attack classifier with richer context.

The key observation is that a model's distribution of maximum confidence scores across many inputs characterizes its confidence behavior. A model that is uniformly confident on all inputs (member-like behavior) differs from one that is confident on some and unsure on others. The \textit{variance} of this distribution captures this spread.

\subsubsection{Feature Construction}

The variance feature is computed as follows:

\begin{enumerate}[leftmargin=*]
    \item \textbf{Pool construction.} During shadow model training, the victim API is queried on each shadow model's training data. The max-confidence scores $\max_c \sigma(z_c)$ from all these queries are pooled into a single vector $\mathbf{v} = [\max(\mathbf{p}_1), \max(\mathbf{p}_2), \ldots]$.
    
    \item \textbf{Bootstrap variance.} For each sample in the attack dataset, $k = 100$ values are drawn uniformly at random (with replacement) from $\mathbf{v}$, and their variance is computed:
    \begin{equation}
        \text{Var}_{\text{bootstrap}} = \text{Var}\big(\{v_{j_1}, v_{j_2}, \ldots, v_{j_k}\}\big), \quad j_i \sim \text{Uniform}(1, |\mathbf{v}|)
    \end{equation}
    This produces slightly different variance values per sample due to the randomness of bootstrapping.
\end{enumerate}

\subsubsection{Applicability}

The variance-of-max feature can be appended to \textit{any} feature vector used for MIA, since it captures a model-level property independent of the specific per-sample features. In our proposed fluctuation-based approach (Section~\ref{sec:fluctuation_mia}), the model-level variance statistic provides the attack classifier with richer context about the victim's overall confidence behavior, helping disambiguate borderline cases where per-sample features alone are insufficient.

In this paper, we evaluate it in the context of the fluctuation-based attack by appending $\text{Var}_{\text{bootstrap}}$ to the 17-dimensional fluctuation feature vector, producing an 18-dimensional feature. Importantly, the variance is computed from the victim's \textit{clean} (non-adversarial) max-confidence scores, ensuring it captures the model's natural confidence behavior rather than its adversarial response.

The shadow models need only be trained \textbf{once} and are shared between the base fluctuation attack and its variance-enhanced version. Only the attack dataset construction and attack model training differ, eliminating redundant computation.

% ─── IV-E. FLUCTUATION-BASED MIA ─────────────────────────────────────────────

\subsection{Fluctuation-Based MIA (Our Approach)}
\label{sec:fluctuation_mia}

We propose a novel \textit{Fluctuation-Based Membership Inference Attack} that directly exploits the robust overfitting phenomenon described in Section~\ref{sec:robust_overfitting}. The core idea is simple yet powerful: instead of training the attack model on raw confidence scores, we train it on the \textit{fluctuation}---the absolute change in confidence scores when adversarial noise is applied to the input.

Crucially, this attack does \textit{not} require the attacker to know whether or how the victim was adversarially trained. The attacker simply probes the victim's response to perturbations as a general-purpose strategy. If the victim was adversarially trained (and thus exhibits robust overfitting), the probing reveals a strong membership signal. If the victim was trained normally, the probing reveals little to no signal, and the attack gracefully degrades to near-random performance without causing harm to the attacker.

% DIAGRAM PLACEHOLDER 5
\begin{figure*}[t]
    \centering
    \includegraphics[width=\textwidth]{images/img5.png}
    \caption{Our proposed Fluctuation-Based MIA pipeline. The key elements are: (1) shadow models are adversarially trained to induce robust overfitting, (2) features are fluctuations $|$clean $-$ adversarial$|$ rather than raw confidences, and (3) victim inference uses transfer-based adversarial examples crafted from shadow model gradients (the victim remains a complete black box).}
    \label{fig:fluctuation_mia}
\end{figure*}

\subsubsection{Key Insight: The Fluctuation Gap}

Define the \textit{fluctuation vector} for an input $x$ as:
\begin{equation}
    \mathbf{z}(x) = \left| f(x) - f(x_{adv}) \right| \in \mathbb{R}^{15}
\end{equation}
where $f(\cdot)$ denotes the model's sigmoid output vector and $x_{adv} = x + \varepsilon \cdot \text{sign}(\nabla_x \mathcal{L})$ is the FGSM adversarial example. Each element $z_c(x)$ measures how much the model's confidence in class $c$ changes under adversarial perturbation.

For a model exhibiting robust overfitting:
\begin{itemize}[leftmargin=*]
    \item \textbf{Members:} The model has been explicitly trained to resist adversarial perturbations on these images. Its outputs are \textit{stable} under noise, yielding \textit{small} fluctuation: $\|\mathbf{z}(x)\| \approx 0.027$.
    \item \textbf{Non-members:} The model's robustness does not generalize. Its outputs are \textit{volatile} under noise, yielding \textit{large} fluctuation: $\|\mathbf{z}(x)\| \approx 0.074$.
\end{itemize}

This $\sim$3$\times$ difference in fluctuation magnitude provides a strong signal for membership inference that is entirely invisible in the clean confidence scores.



\subsubsection{Stage 1: Adversarial Shadow Model Training}

Our approach adversarially trains the shadow models to deliberately induce robust overfitting within each shadow model. This is a strategic choice by the attacker, not one that requires knowledge of the victim's training procedure. The rationale is as follows:

\begin{enumerate}[leftmargin=*]
    \item Adversarial training is known to cause robust overfitting~\cite{rice2020overfitting}, which creates a differential fluctuation between members and non-members.
    \item By adversarially training the shadow models, the attacker ensures that each shadow model \textit{also} exhibits this differential---allowing the attack dataset to capture the fluctuation-based membership signal.
    \item At inference time, the learned relationship ``low fluctuation $\Rightarrow$ member'' transfers to \textit{any} victim model that exhibits robust overfitting, regardless of how the victim was trained.
\end{enumerate}

If the victim does \textit{not} exhibit robust overfitting (e.g., a standard baseline model), then the fluctuation signal will be weak for both members and non-members of the victim, and the attack model will produce near-random predictions. The attack thus gracefully degrades rather than producing false positives.

Each shadow model is trained with FGSM adversarial training using $\varepsilon = 0.01$ and $\alpha = 0.5$ for 40 epochs. The increased epoch count ensures the shadow models overfit sufficiently---they must faithfully replicate the memorization behavior that creates the fluctuation gap.

\subsubsection{Stage 2: Fluctuation Feature Extraction}

For each shadow model's member and non-member subsets, we compute fluctuation features instead of raw confidence scores:

\begin{enumerate}[leftmargin=*]
    \item \textbf{Clean inference.} Query the shadow model on the clean image to obtain $\mathbf{p}_{\text{clean}} = \sigma(f_s(x)) \in [0, 1]^{15}$.
    
    \item \textbf{Adversarial generation.} Generate an FGSM adversarial example using the shadow model's \textit{own} gradients: $x_{adv} = x + \varepsilon \cdot \text{sign}(\nabla_x \mathcal{L}(f_s(x), \hat{y}))$, where $\hat{y}$ are the pseudo-labels obtained by querying the victim API. As discussed in Section~\ref{sec:standard_mia}, these pseudo-labels ensure the shadow model's gradients are aligned with the victim's decision boundary, making the perturbation directions transferable.
    
    \item \textbf{Adversarial inference.} Query the shadow model on the adversarial image to obtain $\mathbf{p}_{\text{adv}} = \sigma(f_s(x_{adv}))$.
    
    \item \textbf{Fluctuation computation.} Compute the per-class absolute difference:
    \begin{equation}
        \mathbf{z} = |\mathbf{p}_{\text{clean}} - \mathbf{p}_{\text{adv}}| \in \mathbb{R}^{15}
    \end{equation}
\end{enumerate}

Each sample in the attack dataset consists of 17 features:
\begin{equation}
    \text{features} = \big(\underbrace{z_1, \ldots, z_{15}}_{\text{fluctuation}},\; \underbrace{\max(\mathbf{p}_{\text{clean}})}_{\text{max conf (clean)}},\; \underbrace{\arg\max(\mathbf{p}_{\text{clean}})}_{\text{max class (clean)}}\big)
\end{equation}
Note that the max confidence and max class are computed from the \textit{clean} output, not the adversarial output.

\subsubsection{Stage 3: Attack Model Training}

A Gradient Boosting classifier is trained on the 17-dimensional fluctuation feature vectors to predict membership. Because the fluctuation gap between members and non-members is substantially larger than the clean confidence gap, the attack model can learn a highly discriminative decision boundary from these features.

\subsubsection{Stage 4: Transfer-Based Victim Inference}

At inference time, the attacker must generate adversarial examples to probe the victim model. However, the victim is a \textit{black-box}---the attacker cannot compute gradients through the victim. We solve this using a \textit{transfer attack}~\cite{papernot2017practical}: adversarial examples crafted using a shadow model's gradients are applied to the victim model. This is possible because the shadow models were trained to mimic the victim's decision boundary, so their gradient directions are correlated.

\begin{algorithm}[t]
\caption{Fluctuation-Based MIA -- Inference}
\label{alg:fluctuation_mia}
\begin{algorithmic}[1]
\Require Victim API $f_V$ (black-box), trained shadow model $f_S$, eval image $x$, perturbation $\varepsilon$, trained attack model $\mathcal{A}$
\State $\mathbf{p}_{\text{clean}} \gets \sigma(f_V(x))$ \Comment{Query victim: clean image}
\State $\hat{y} \gets \mathbb{1}[\mathbf{p}_{\text{clean}} > 0.5]$ \Comment{Pseudo-labels from victim}
\State $g \gets \text{sign}\big(\nabla_x \mathcal{L}(f_S(x), \hat{y})\big)$ \Comment{Gradient from \textit{shadow}}
\State $x_{adv} \gets \text{clamp}(x + \varepsilon \cdot g,\; -3, 3)$ \Comment{FGSM via shadow}
\State $\mathbf{p}_{\text{adv}} \gets \sigma(f_V(x_{adv}))$ \Comment{Query victim: adversarial}
\State $\mathbf{z} \gets |\mathbf{p}_{\text{clean}} - \mathbf{p}_{\text{adv}}|$ \Comment{Fluctuation vector}
\State $\text{feat} \gets (\mathbf{z},\; \max(\mathbf{p}_{\text{clean}}),\; \arg\max(\mathbf{p}_{\text{clean}}))$
\State \Return $\mathcal{A}(\text{feat})$ \Comment{Predict: member or non-member}
\end{algorithmic}
\end{algorithm}

The transfer attack works because the shadow models were trained to mimic the victim's decision boundary. The gradient directions of the shadow model are correlated with those of the victim, so adversarial perturbations that change the shadow model's output also change the victim's output in a similar manner. Crucially, the \textit{magnitude} of the change differs between members and non-members of the \textit{victim's} training set---this is the signal the attack model has learned to detect.

\subsubsection{Why This Approach Works}

The fluctuation-based attack succeeds because it directly probes the mechanism of robust overfitting. A confidence-based approach examines the model's output in a ``resting state'' (clean input), where the adversarially trained model has learned to produce uniform outputs that reveal nothing. Our approach actively \textit{perturbs} the model and measures its response---analogous to tapping a structure and listening for resonance patterns that reveal its internal construction. Members produce a ``solid'' response (low fluctuation); non-members produce a ``hollow'' response (high fluctuation).

This analogy extends to why the attack requires no knowledge of the victim's training: a doctor tapping a patient's chest does not need to know the patient's medical history to interpret the sound. Similarly, the fluctuation probe is a general diagnostic technique whose results are informative whenever the model exhibits any form of data-dependent robustness.


% ═══════════════════════════════════════════════════════════════════════════════
% V. PERFORMANCE METRICS AND RESULTS
% ═══════════════════════════════════════════════════════════════════════════════

\section{Performance Metrics and Results}
\label{sec:results}

\subsection{Experimental Setup}

All experiments are conducted on a single NVIDIA GeForce RTX 4090 GPU with 24 GB VRAM. The victim models are trained using PyTorch with mixed-precision training (FP16). Shadow models are trained on the same hardware. The MIA evaluation uses 5,000 held-out member images and 5,000 held-out non-member images that are disjoint from the attacker's pool.

\subsection{Evaluation Metrics}

\subsubsection{Evasion Attack}
We measure the macro-average Area Under the ROC Curve (AUC) under FGSM perturbation at varying $\varepsilon$ values. A lower AUC indicates a more successful evasion attack. We report:
\begin{itemize}[leftmargin=*]
    \item \textbf{Clean AUC:} Model performance with no perturbation ($\varepsilon = 0$).
    \item \textbf{Adversarial AUC:} Model performance under FGSM at each $\varepsilon$.
    \item \textbf{Robustness gain:} Difference in adversarial AUC between the adversarial and baseline models.
\end{itemize}

\subsubsection{Membership Inference Attack}
We evaluate the MIA using standard binary classification metrics on the balanced evaluation set:
\begin{itemize}[leftmargin=*]
    \item \textbf{Accuracy:} Fraction of correct membership predictions. Random guessing yields 50\%.
    \item \textbf{Precision:} Fraction of predicted members that are true members.
    \item \textbf{Recall:} Fraction of true members correctly identified.
    \item \textbf{F1 Score:} Harmonic mean of precision and recall.
\end{itemize}



\subsection{Evasion Attack Results}

Figure~\ref{fig:epsilon_vs_fliprate} shows the prediction flip rate as a function of FGSM perturbation magnitude $\varepsilon$ for both victim models. Table~\ref{tab:evasion} provides the exact values.

\begin{figure}[t]
    \centering
    \includegraphics[width=\columnwidth]{images/graph1.png}
    \caption{Prediction flip rate under FGSM evasion attack at increasing $\varepsilon$. The baseline model (blue) suffers rapidly increasing flip rates up to 26.1\% at $\varepsilon=0.02$, while the adversarially trained model (orange) keeps flip rates below 5.4\% across all perturbation magnitudes, demonstrating strong evasion defense.}
    \label{fig:epsilon_vs_fliprate}
\end{figure}


\begin{table}[!ht]
    \centering
    \caption{Evasion Attack: Classification AUC under FGSM at Varying $\varepsilon$}
    \label{tab:evasion}
    \begin{tabular}{lccccc}
        \toprule
        \textbf{$\varepsilon$} & \textbf{Baseline AUC} & \textbf{Adversarial AUC} & \textbf{Baseline Flip} & \textbf{Adversarial Flip} \\
        \midrule
        0.000 (clean) & 0.8278 & 0.7516 & --- & --- \\
        0.001         & 0.6832 & 0.5115 & 7.4\% & 4.2\% \\
        0.002         & 0.5842 & 0.4669 & 12.9\% & 5.0\% \\
        0.005         & 0.4396 & 0.7935 & 20.5\% & 1.9\% \\
        0.010         & 0.3591 & 0.9044 & 24.2\% & 5.4\% \\
        0.020         & 0.3110 & 0.8264 & 26.1\% & 3.9\% \\
        \bottomrule
    \end{tabular}
\end{table}

\subsection{Membership Inference Attack Results}

Tables~\ref{tab:mia_baseline} and~\ref{tab:mia_adversarial} present the detailed fluctuation-based MIA results for each victim model across all four attack classifiers and both feature strategies.

\begin{table}[t]
    \centering
    \caption{MIA Results --- Baseline Victim Model}
    \label{tab:mia_baseline}
    \resizebox{\columnwidth}{!}{%
    \begin{tabular}{llcccc}
        \toprule
        \textbf{Feature Strategy} & \textbf{Attack Classifier} & \textbf{Acc} & \textbf{Prec} & \textbf{Rec} & \textbf{F1} \\
        \midrule
        \multirow{4}{*}{Fluctuation}
            & Gradient Boosting    & 0.5016 & 0.5256 & 0.0328 & 0.0617 \\
            & Random Forest        & 0.5100 & 0.5144 & 0.3580 & 0.4222 \\
            & MLP                  & 0.5047 & 0.5050 & 0.4714 & 0.4876 \\
            & Logistic Regression  & 0.5042 & 0.5265 & 0.0834 & 0.1440 \\
        \midrule
        \multirow{4}{*}{\shortstack[l]{Fluctuation\\+ Variance}}
            & Gradient Boosting    & 0.5004 & 0.5071 & 0.0286 & 0.0541 \\
            & Random Forest        & 0.4981 & 0.4959 & 0.2278 & 0.3122 \\
            & MLP                  & 0.5055 & 0.5051 & 0.5490 & 0.5261 \\
            & Logistic Regression  & 0.5043 & 0.5280 & 0.0812 & 0.1408 \\
        \bottomrule
    \end{tabular}%
    }
    \vspace{2pt}
    {\footnotesize Random baseline: 0.5000. All classifiers perform near random guessing, confirming the baseline model leaks no membership signal.}
\end{table}

\begin{table}[t]
    \centering
    \caption{MIA Results --- Adversarial Victim Model ($\varepsilon=0.02$)}
    \label{tab:mia_adversarial}
    \resizebox{\columnwidth}{!}{%
    \begin{tabular}{llcccc}
        \toprule
        \textbf{Feature Strategy} & \textbf{Attack Classifier} & \textbf{Acc} & \textbf{Prec} & \textbf{Rec} & \textbf{F1} \\
        \midrule
        \multirow{4}{*}{Fluctuation}
            & Gradient Boosting    & \textbf{0.6730} & 0.6048 & 0.9984 & \textbf{0.7533} \\
            & Random Forest        & 0.6618 & 0.6227 & 0.8212 & 0.7083 \\
            & MLP                  & 0.6128 & 0.5636 & 0.9996 & 0.7208 \\
            & Logistic Regression  & 0.5938 & 0.5518 & 1.0000 & 0.7111 \\
        \midrule
        \multirow{4}{*}{\shortstack[l]{Fluctuation\\+ Variance}}
            & Gradient Boosting    & \textbf{0.6744} & 0.6059 & 0.9978 & \textbf{0.7540} \\
            & Random Forest        & 0.6641 & 0.6230 & 0.8314 & 0.7122 \\
            & MLP                  & 0.6038 & 0.5579 & 0.9996 & 0.7161 \\
            & Logistic Regression  & 0.5936 & 0.5516 & 1.0000 & 0.7110 \\
        \bottomrule
    \end{tabular}%
    }
    \vspace{2pt}
    {\footnotesize Gradient Boosting with Fluctuation+Variance achieves the best accuracy (67.44\%) and F1 (0.754), confirming the adversarial model is vulnerable to membership inference.}
\end{table}

\subsection{The Robustness--Privacy Tradeoff}

Figure~\ref{fig:tradeoff} visualizes the core finding: the tradeoff between evasion attack defense and MIA vulnerability.

\begin{figure}[t]
    \centering
    \includegraphics[width=\columnwidth]{images/graph2.png}
    \caption{The robustness--privacy tradeoff. \textbf{Left:} The baseline model suffers a 26.1\% flip rate under FGSM ($\varepsilon=0.02$), while the adversarial model reduces this to just 3.9\%, demonstrating strong evasion defense. \textbf{Right:} This robustness comes at a privacy cost---the adversarial model is vulnerable to fluctuation-based MIA (67.4\% accuracy) while the baseline remains at random-guessing level (50.2\%). The dashed line indicates random guessing.}
    \label{fig:tradeoff}
\end{figure}




\subsection{Discussion}

The results reveal several important findings:

\textbf{1. Adversarial training effectively defends against evasion.} The adversarially trained model maintains consistently low prediction flip rates ($\leq$5.4\%) under FGSM perturbation across all $\varepsilon$ values (Fig.~\ref{fig:epsilon_vs_fliprate}), while the baseline model suffers flip rates up to 26.1\%, validating adversarial training as an effective evasion defense.

\textbf{2. Fluctuation-based MIA exploits robust overfitting.} Our proposed attack achieves $\sim$70\% accuracy on the adversarially trained model---a 20 percentage point improvement over random guessing---while remaining at $\sim$51\% on the baseline. This confirms that robust overfitting, not standard overfitting, is the vulnerability, and that the fluctuation probing technique successfully detects it without requiring knowledge of the victim's training procedure.

\textbf{3. The robustness--privacy tradeoff is real.} Taken together, findings 1 and 2 establish a clear inverse relationship: the adversarial model gains evasion robustness but loses privacy; the baseline model is evasion-vulnerable but privacy-safe. Defending against one attack class opens vulnerability to the other.

\textbf{4. Variance enhancement improves attack accuracy.} The addition of the $\text{Var}_{\text{bootstrap}}$ feature provides a meaningful boost to the fluctuation-based MIA, complementing the per-sample fluctuation features with model-level behavioral context. As discussed in Section~\ref{sec:variance_mia}, this enhancement is general-purpose and could similarly improve confidence-based MIA approaches by enriching the feature space with distributional information.



% ═══════════════════════════════════════════════════════════════════════════════
% VI. CONCLUSION
% ═══════════════════════════════════════════════════════════════════════════════

\section{Conclusion}
\label{sec:conclusion}

We have demonstrated a fundamental \textit{robustness--privacy tradeoff} in adversarially trained medical image classifiers. While FGSM adversarial training effectively defends against evasion attacks, it simultaneously introduces robust overfitting---a phenomenon where the model's adversarial robustness is specific to its training data and fails to generalize. Our novel fluctuation-based MIA exploits this asymmetry by measuring the differential stability of model outputs under adversarial perturbation, achieving $\sim$70\% attack accuracy on the adversarially trained model compared to $\sim$51\% on the standard baseline---all without requiring any knowledge of the victim's training procedure.

The baseline model, which was not adversarially trained, exhibits uniformly high fluctuation for both members and non-members, confirming that robust overfitting---not standard overfitting---is the vulnerability exploited by our attack. Additionally, the variance-of-max feature enhancement we introduce is a general-purpose technique that improves attack accuracy across different MIA strategies.

These findings underscore a critical challenge for medical AI deployment: defending against one class of adversarial threats (evasion) can inadvertently create vulnerability to another (privacy). Future work should explore mitigation strategies such as:

\begin{enumerate}[leftmargin=*]
    \item \textbf{Differential privacy.} Adding calibrated noise to gradients during training can provably limit membership inference risk, but may reduce model utility.
    \item \textbf{Robust overfitting mitigation.} Techniques such as early stopping on adversarial validation AUC~\cite{rice2020overfitting}, adversarial weight perturbation~\cite{wu2020adversarial}, or stochastic weight averaging~\cite{chen2021robust} may reduce the memorization gap while preserving robustness.
    \item \textbf{Prediction perturbation.} Adding noise to the model's output confidence scores before returning them to the API can mask the fluctuation signal, though this degrades diagnostic utility.
\end{enumerate}

Addressing the robustness--privacy tradeoff is essential for building trustworthy medical AI systems that are simultaneously secure against adversarial manipulation and protective of patient privacy.


% ═══════════════════════════════════════════════════════════════════════════════
% REFERENCES
% ═══════════════════════════════════════════════════════════════════════════════

\begin{thebibliography}{00}

\bibitem{goodfellow2015explaining}
I.~J.~Goodfellow, J.~Shlens, and C.~Szegedy, ``Explaining and harnessing adversarial examples,'' in \textit{Proc. Int. Conf. Learning Representations (ICLR)}, 2015.

\bibitem{madry2018towards}
A.~Madry, A.~Makelov, L.~Schmidt, D.~Tsipras, and A.~Vladu, ``Towards deep learning models resistant to adversarial attacks,'' in \textit{Proc. Int. Conf. Learning Representations (ICLR)}, 2018.

\bibitem{shokri2017membership}
R.~Shokri, M.~Stronati, C.~Song, and V.~Shmatikov, ``Membership inference attacks against machine learning models,'' in \textit{Proc. IEEE Symp. Security and Privacy (S\&P)}, 2017, pp.~3--18.

\bibitem{song2019privacy}
L.~Song, R.~Shokri, and P.~Mittal, ``Privacy risks of securing machine learning models against adversarial examples,'' in \textit{Proc. ACM Conf. Computer and Communications Security (CCS)}, 2019.

\bibitem{rice2020overfitting}
L.~Rice, E.~Wong, and J.~Z.~Kolter, ``Overfitting in adversarially robust deep learning,'' in \textit{Proc. Int. Conf. Machine Learning (ICML)}, 2020.

\bibitem{rajpurkar2017chexnet}
P.~Rajpurkar \textit{et al.}, ``CheXNet: Radiologist-level pneumonia detection on chest X-rays with deep learning,'' \textit{arXiv preprint arXiv:1711.05225}, 2017.

\bibitem{huang2017densely}
G.~Huang, Z.~Liu, L.~Van~Der~Maaten, and K.~Q.~Weinberger, ``Densely connected convolutional networks,'' in \textit{Proc. IEEE Conf. Computer Vision and Pattern Recognition (CVPR)}, 2017.

\bibitem{wang2017chestxray}
X.~Wang \textit{et al.}, ``ChestX-Ray8: Hospital-scale chest X-ray database and benchmarks on weakly-supervised classification and localization of common thorax diseases,'' in \textit{Proc. IEEE Conf. Computer Vision and Pattern Recognition (CVPR)}, 2017.

\bibitem{salem2019ml}
A.~Salem, Y.~Zhang, M.~Humbert, P.~Berrang, M.~Fritz, and M.~Backes, ``ML-Leaks: Model and data independent membership inference attacks and defenses on machine learning models,'' in \textit{Proc. Network and Distributed System Security Symposium (NDSS)}, 2019.

\bibitem{szegedy2014intriguing}
C.~Szegedy \textit{et al.}, ``Intriguing properties of neural networks,'' in \textit{Proc. Int. Conf. Learning Representations (ICLR)}, 2014.

\bibitem{carlini2017towards}
N.~Carlini and D.~Wagner, ``Towards evaluating the robustness of neural networks,'' in \textit{Proc. IEEE Symp. Security and Privacy (S\&P)}, 2017.

\bibitem{deng2009imagenet}
J.~Deng, W.~Dong, R.~Socher, L.-J.~Li, K.~Li, and L.~Fei-Fei, ``ImageNet: A large-scale hierarchical image database,'' in \textit{Proc. IEEE Conf. Computer Vision and Pattern Recognition (CVPR)}, 2009.

\bibitem{gulshan2016development}
V.~Gulshan \textit{et al.}, ``Development and validation of a deep learning algorithm for detection of diabetic retinopathy in retinal fundus photographs,'' \textit{JAMA}, vol.~316, no.~22, pp.~2402--2410, 2016.

\bibitem{esteva2017dermatologist}
A.~Esteva \textit{et al.}, ``Dermatologist-level classification of skin cancer with deep neural networks,'' \textit{Nature}, vol.~542, no.~7639, pp.~115--118, 2017.

\bibitem{finlayson2019adversarial}
S.~G.~Finlayson, J.~D.~Bowers, J.~Ito, J.~L.~Zittrain, A.~L.~Beam, and I.~S.~Kohane, ``Adversarial attacks on medical machine learning,'' \textit{Science}, vol.~363, no.~6433, pp.~1287--1289, 2019.

\bibitem{yeom2018privacy}
S.~Yeom, I.~Giacomelli, M.~Fredrikson, and S.~Jha, ``Privacy risk in machine learning: Analyzing the connection to overfitting,'' in \textit{Proc. IEEE Computer Security Foundations Symposium (CSF)}, 2018.

\bibitem{song2021systematic}
L.~Song and P.~Mittal, ``Systematic evaluation of privacy risks of machine learning models,'' in \textit{Proc. USENIX Security Symposium}, 2021.

\bibitem{wu2020adversarial}
D.~Wu, S.~Xia, and Y.~Wang, ``Adversarial weight perturbation helps robust generalization,'' in \textit{Proc. Advances in Neural Information Processing Systems (NeurIPS)}, 2020.

\bibitem{paschali2018generalizability}
M.~Paschali, S.~Conjeti, F.~Navarro, and N.~Navab, ``Generalizability vs. robustness: Investigating medical imaging networks using adversarial examples,'' in \textit{Proc. MICCAI}, 2018.

\bibitem{papernot2017practical}
N.~Papernot, P.~McDaniel, I.~Goodfellow, S.~Jha, Z.~B.~Celik, and A.~Swami, ``Practical black-box attacks against machine learning,'' in \textit{Proc. ACM Asia CCS}, 2017.

\bibitem{guo2018countering}
C.~Guo, M.~Rana, M.~Cissé, and L.~van der Maaten, ``Countering adversarial images using input transformations,'' in \textit{Proc. ICLR}, 2018.

\bibitem{papernot2016distillation}
N.~Papernot, P.~McDaniel, X.~Wu, S.~Jha, and A.~Swami, ``Distillation as a defense to adversarial perturbations against deep neural networks,'' in \textit{Proc. IEEE S\&P}, 2016.

\bibitem{cohen2019certified}
J.~Cohen, E.~Rosenfeld, and J.~Z.~Kolter, ``Certified adversarial robustness via randomized smoothing,'' in \textit{Proc. ICML}, 2019.

\bibitem{choquette2021label}
C.~A.~Choquette-Choo, F.~Tramer, N.~Carlini, and N.~Papernot, ``Label-only membership inference attacks,'' in \textit{Proc. ICML}, 2021.

\bibitem{carlini2022membership}
N.~Carlini \textit{et al.}, ``Membership inference attacks from first principles,'' in \textit{Proc. IEEE S\&P}, 2022.

\bibitem{chen2021robust}
T.~Chen, Z.~Zhang, S.~Liu, S.~Chang, and Z.~Wang, ``Robust overfitting may be mitigated by properly learned smoothing,'' in \textit{Proc. ICLR}, 2021.

\end{thebibliography}

\end{document}
